\begin{abstract}
Graph neural networks are widely proposed for detecting illicit Bitcoin
transactions, yet under strict evaluation they repeatedly lose to tree
ensembles, leaving the motivation for graph learning unclear. We argue that
prior work used the wrong graph and test this with a time-sliced
transaction--address knowledge graph built from Elliptic++, in which addresses
are instantiated per time step, carry only causally computed history, and no
edge crosses a time step, so that evaluation is strictly inductive by
construction. Guided by a recent reverse engineering of the dataset, we also
remove features that are computed after the test period. Three findings
emerge. First, these leaked features inflate the illicit-class F1 of standard
baselines from \XgbCleanFone{} to \XgbAllFone{}. Second, the knowledge graph exposes
entity structure that is far more homophilous than the transaction graph
(class-adjusted homophily 0.92 versus 0.38 on the test period), and causal
knowledge-graph features raise a leakage-controlled XGBoost from F1
\XgbCleanFone{} to \XgbCleanKgFone{}. Third, relational message passing over the same graph achieves the best leakage-controlled
validation score and improves over other neural models on previously unseen
entities, but raises many false alarms on recurring licit entities and is the
weakest model on the full test period.
Graph structure helps when it is distilled into causal features for robust
learners, not when it is learned end to end.
\end{abstract}
